\documentclass[10pt,a4paper]{article}
\usepackage[utf8]{inputenc}
\usepackage{booktabs}
\usepackage[a4paper,left=3cm,right=2cm]{geometry}
\usepackage{amsmath}
\usepackage[thmmarks, amsmath, thref]{ntheorem}
\usepackage{amsfonts}
\usepackage{amssymb}
\usepackage{fancyhdr}
\pagestyle{fancy}
\usepackage{anysize}
\usepackage{graphicx}
\usepackage{float}
\usepackage{hyperref}
\usepackage{enumitem}
\usepackage{verbatim}
\newtheorem{theorem}{Theorem}
\theoremsymbol{\ensuremath{\blacksquare}}
\newtheorem*{solution}{Solución:}

\usepackage{xcolor}
\usepackage{listings}
% Margenes y formalidades

\textwidth 6.25in % Margen Derecho y ancho del texto
\textheight 8.5in
\oddsidemargin 0in % Margen Izquierdo
\headheight 0in % Largo del texto
\leftmargin 1in
\parindent 0em % Salto entre la indentación (Sangría)https://www.overleaf.com/4519127669sdpptkhvsbfn
\topmargin -0.5in % Margen Superior
\parskip 0.5ex % Salto entre parrafos
\baselineskip 1pt



\lhead{
\includegraphics[scale=0.08]{Figuras/logoUSM.png}} % texto izquierda de la cabecera
%\chead{}                                      % texto centro de la cabecera
\rhead{Universidad Técnica Federico Santa María\\Ingeniería Comercial } % número de página a la derecha
\renewcommand{\headrulewidth}{0.4pt} % grosor de la línea de la cabecera
\renewcommand{\footrulewidth}{0.4pt} % grosor de la línea del pie

\begin{document}
\vspace{5cm}



\begin{center}
\begin{large}
\textbf{\\CERTAMEN 2:}\textbf{ ECONOMETRÍA ICS-013}
\end{large}\\
\vspace{0.2cm}
II Semestre 2025\\
Profesores: Jocelyn Tapia Stefanoni, Franco Antonucci Fernández y Ricardo Pruneda Urz\'ua\\
%Ayudante: 
\end{center}

\textbf{Instrucciones:} 

\begin{itemize}

    \item Usted dispone de 90 minutos  para intentar los 100 puntos del certamen.
    \item Cuide la presentación y redacción de sus preguntas.
    \item No puede utilizar celular, computador o tablet, pero sí una calculadora TI.
     
\end{itemize}

\section{Comentes y preguntas cortas (40 puntos)}

Responda solo 4 de las siguientes preguntas.
\begin{enumerate}
    \item \textbf{RdeA9.1} (10 puntos) 
       Muestre que si se cumplen los supuestos de Teorema de Gauss-Markov (los primeros 4) el estimador de MCO $ \hat{\beta} = (X^{\top}X)^{-1}X^{\top}y$
es insesgado. Señale como utilizó los supuestos en su demostración.

 \item \textbf{RdeA9.1} (10 puntos) 
Suponga un modelo dado por:

$$\text{PAES}_i = \beta_0 + \beta_1 \cdot expend_i + \beta_2 \cdot ingreso_i + u$$

Donde PAES es el puntaje de ingreso que un alumno obtuvo en la prueba de acceso a la educación superior, \textit{expend} es el gasto en educación en el estudiante e \textit{ingreso} es el ingreso familiar promedio del estudiante. Se esperaría que \textit{expend} e \textit{ingreso} estén correlacionadas, puesto que las comunas donde en los que el ingreso familiar promedio es mayor tienden a gastar más en educación por estudiante. Es por esto que el modelo no cumple con el supuesto de colinealidad y, por lo tanto, las estimaciones serán insesgadas e inconsistentes. Comente.

 \item \textbf{RdeA9.1} (10 puntos) 
En la tabla \ref{tabla_reg}  se muestran los resultados de una investigación que busca identificar los determinantes de los salarios de los profesores (avgsal).  Se utilizan las siguientes variables explicativas 

\begin{itemize}
   \item \texttt{lunch}: porcentaje de estudiantes elegibles para almuerzo gratuito.
  \item \texttt{enrol}: matrícula (número de estudiantes).
  \item \texttt{staff}: personal por cada 1000 estudiantes.
 
  \item \texttt{math4}: porcentaje que aprueba la prueba de matemáticas de 4.\textsuperscript{o} grado.
  \item \texttt{story4}: porcentaje que aprueba la prueba de lectura de 4.\textsuperscript{o} grado.
  \item \texttt{bs}: razón beneficios/salario, definida como \texttt{avgben/avgsal}.
\item  \texttt{lenrol}: logaritmo natural de \texttt{enrol}.
  \item \texttt{lstaff}: logaritmo natural de \texttt{staff}.
\end{itemize}

\begin{table}[H] \centering 
  \caption{Salarios de los profesores }  \label{tabla_reg}%\caption{Comparación: Regresiones} 
\begin{tabular}[H]{@{\extracolsep{5pt}}lcc} 
\\[-1.8ex]\hline 
\hline \\[-1.8ex] 
 & \multicolumn{2}{c}{\textit{Dependent variable:}} \\ 
\cline{2-3} 
\\[-1.8ex] & \multicolumn{2}{c}{log(avgsal)} \\ 
\\[-1.8ex] & (1) & (2)\\ 
\hline \\[-1.8ex] 
 bs & $-$0.516$^{***}$ & $-$0.518$^{***}$ \\ 
  & (0.110) & (0.110) \\ 
  & & \\ 
 lenrol & $-$0.028$^{***}$ & $-$0.028$^{***}$ \\ 
  & (0.008) & (0.008) \\ 
  & & \\ 
 lstaff & $-$0.691$^{***}$ & $-$0.688$^{***}$ \\ 
  & (0.018) & (0.018) \\ 
  & & \\ 
 lunch & $-$0.001$^{***}$ & $-$0.001$^{***}$ \\ 
  & (0.0002) & (0.0002) \\ 
  & & \\ 
 math4 &  & 0.0004 \\ 
  &  & (0.0004) \\ 
  & & \\ 
 story4 &  & $$0.001 \\ 
  &  & (0.0005) \\ 
  & & \\ 
 Constant & 13.831$^{***}$ & 13.850$^{***}$ \\ 
  & (0.110) & (0.113) \\ 
  & & \\ 
\hline \\[-1.8ex] 
Observations & 1,848 & 1,848 \\ 
R$^{2}$ & 0.488 & 0.493 \\ 

\hline 
\hline \\[-1.8ex] 
\textit{Note:}  & \multicolumn{2}{r}{$^{*}$p$<$0.1; $^{**}$p$<$0.05; $^{***}$p$<$0.01} \\ 
\end{tabular}
\end{table} 

Debiese incluir las variables \texttt{math4} y \texttt{story4} en la investigación? Sustente su respuesta mediante la formulación y explicación de un test de hipótesis que permita evaluar su relevancia estadística en el modelo ($
F_{0.05, (2, 1841)} \approx 3.00$).
 
 \item \textbf{RdeA9.1} (10 puntos)  
{Considere el siguiente modelo del tiempo de entrega en función de la carga de trabajo:}

\[
\text{tiempo} = 12 - 0.7\,\text{carga} + 0.02\,\text{carga}^2
\]

\text{donde el tiempo se mide en horas y la carga representa el número de pedidos procesados por día.}

\begin{enumerate}
    \item (5 puntos) ¿Cuál es la interpretación de $\beta_1$ y $\beta_2$?

    \item (5 puntos) Determine el nivel de carga que minimiza el tiempo de entrega.
\end{enumerate}

 
 \item \textbf{RdeA9.1} (10 puntos)  
{Suponga un modelo aplicado a las ventas dado por:}

\[
\ln(\mathrm{ventas}) = \beta_0 + \beta_1\,\mathrm{publicidad} + \beta_2\,\mathrm{precio} + u
\]

\noindent
\textbf{Definiciones y unidades:}
\begin{itemize}
    \item \textbf{ventas}: número de unidades vendidas (logaritmo natural)
    \item \textbf{publicidad}: inversión en publicidad (miles de dólares)
    \item \textbf{precio}: precio del producto (dólares por unidad)
\end{itemize}

\begin{center}
\begin{tabular}{lcc}
\toprule
\textbf{Variable} & \textbf{Coeficiente} & \textbf{Error Estándar} \\
\midrule
Intercepto   & 2.00***  & 0.19 \\
publicidad   & 0.058**  & 0.021 \\
precio       & -0.016*  & 0.008 \\
\bottomrule
\end{tabular}

\vspace{0.5em}

$R^2 = 0.42$ \quad $F = 28.5$ (gl = 2, 92) \\[0.3em]
\textit{Nota:} *p $<$ 0.1; **p $<$ 0.05; ***p $<$ 0.01
\end{center}

 a) (5 puntos) 
Realice una prueba de significancia individual para el parámetro $\beta_{1}$ asociado a la variable \texttt{publicidad}, considerando un nivel de significancia del 10\%. Explique sus resultados.

%Realice la inferencia para el parámetro $\beta_1$ (publicidad) al 10\% de significancia.
% Formule las hipótesis nula y alternativa, calcule el estadístico t correspondiente, indique el valor crítico
% y concluya si el parámetro es significativamente distinto de cero.

b) (5 puntos) Un investigador sostiene que el efecto del precio sobre las ventas es mayor que cero
al 1\% de significancia. ¿Está usted de acuerdo? Justifique su respuesta utilizando los resultados del modelo.
 
 \item \textbf{RdeA9.1} (10 puntos)  Se plantea el modelo:

$$votosA = \beta_0 + \beta_1 \cdot ln(\text{gasto A}) + \beta_2 \cdot ln(\text{gasto B}) + \beta_3 \cdot forpart + u$$

\begin{table}[H]
\centering
\caption{Resultados de la estimación del modelo de regresión}
\label{tab:regresion}
\begin{tabular}{lcccc}
\toprule
\textbf{Variable} & \textbf{Coeficiente} & \textbf{Error estándar} & \textbf{t-valor} & \textbf{Pr(>|t|)} \\
\midrule
Intercepto & 45.07893 & 3.92631 & 11.48 & $<$2e-16$^{***}$ \\
$\ln(\text{gasto A})$ & 6.08332 & 0.38215 & 15.92 & $<$2e-16$^{***}$ \\
$\ln(\text{gasto B})$ & -6.61542 & 0.37882 & -17.46 & $<$2e-16$^{***}$ \\
\textit{forpart} & 0.15196 & 0.06202 & 2.45 & 0.0153$^{*}$ \\
\midrule
\multicolumn{5}{l}{\textit{Signif. codes:} $^{***}p<0.001$, $^{**}p<0.01$, $^{*}p<0.05$} \\
\midrule
\multicolumn{5}{l}{Error estándar residual: 7.712 (169 g.l.)} \\
\multicolumn{5}{l}{$R^2$ = 0.7926 \quad $R^2_{ajustado}$ = 0.7889} \\
\multicolumn{5}{l}{Estadístico F = 215.2 (3; 169) \quad p-valor $<$ 2.2e-16} \\
\bottomrule
\end{tabular}
\end{table}
% \begin{lstlisting}
%  Coefficients:
%             Estimate Std. Error t value Pr(>|t|)    
% (Intercept) 45.07893    3.92631   11.48   <2e-16 ***
% ln gasto A     6.08332    0.38215   15.92   <2e-16 ***
% ln gasto B    -6.61542    0.37882  -17.46   <2e-16 ***
% forpart    0.15196    0.06202    2.45   0.0153 *  
% ---
% Signif. codes:  0 ‘***’ 0.001 ‘**’ 0.01 ‘*’ 0.05 ‘.’ 0.1 ‘ ’ 1

% Residual standard error: 7.712 on 169 degrees of freedom
% Multiple R-squared:  0.7926,	Adjusted R-squared:  0.7889 
% F-statistic: 215.2 on 3 and 169 DF,  p-value: < 2.2e-16
% \end{lstlisting}



Donde $ln(\text{gasto A})$ y $ln(\text{gasto B})$ son el logaritmo natural del gasto en campaña del candidato A y B, respectivamente, $forpart$ es una medida de la fortaleza del partido del candidato A (porcentaje de votos que obtuvo el partido A en la elección presidencial más reciente), y votosA el porcentaje de votos recibidos por el candidato A.\\

Además, se tiene que la matriz de varianzas-covarianzas es:

\begin{table}[H]
    \centering
    \begin{tabular}{c|c|c|c|c}
 & (Intercept) & gasto A & gasto B	& forpart\\\hline
(Intercept)	& 15.4158713 & -0.394937598	& -0.874063434 & -0.176167017\\
gasto A	& -0.3949376 & 0.146038608 & -0.002681489 & -0.006546373\\
gasto B	& -0.8740634 & -0.002681489	& 0.143504821 & 0.003577341\\
forpart	& -0.1761670 & -0.006546373	& 0.003577341 & 0.003846245\\
    \end{tabular}
    %\caption{Caption}
    %\label{tab:placeholder}
\end{table}

Plantee el test de hipótesis de que un aumento de un 1\% en los gastos de A es compensado por un aumento de 1\% en los gastos de B, determine la veracidad de lo anterior al 5\% y explique.\\

    
    \end{enumerate}

\section{Ejercicios (60 puntos)}

Responda solo 2 de las siguientes preguntas.

\begin{enumerate}
    \item \textbf{RdeA9.1} (30 puntos) 
Considere el siguiente modelo de regresión lineal con los supuestos estándares:

\[
Y_i = \beta_0 + \beta_1 X_{1i} + \beta_2 X_{2i} + u_i
\]

Se dispone de la siguiente información:

\[
X'X = 
\begin{pmatrix}
99 & 0 & 0 \\
0 & 120 & 60 \\
0 & 60 & 180
\end{pmatrix},
\qquad
X'Y = 
\begin{pmatrix}
396 \\
72 \\
276
\end{pmatrix}
\]

\[
\sum_{i=1}^{N} Y_i^2 = 2034, 
\qquad 
\sum_{i=1}^{N} (Y_i - \bar{Y})^2 = 450
\]

\begin{enumerate}
    \item (3 puntos) ¿Cuál es el tamaño de la muestra?\\

    \item (13 puntos) Calcule los estimadores de MCO del modelo anterior.\\
    \textbf{Hint:} Si \( Z = \begin{pmatrix} A &{\bf 0} \\ {\bf 0} & B \end{pmatrix} \) con \( A \) y \( B \) matrices cuadradas, entonces
    \[
    Z^{-1} = \begin{pmatrix} A^{-1} & {\bf 0} \\ {\bf 0} & B^{-1} \end{pmatrix}.
    \]
    
    \item (4 puntos) Indique la ecuación del modelo predicho e interprete.\\
    
    \item (5 puntos) ¿Qué debería suceder con $R^2$ si agregamos más variables independientes a la medición?\\


    \item (5 puntos) ¿Qué valor alcanza el $R^2$?\footnote{Hint: Recuerde que $SSE = \sum_{i=1}^N (\hat{Y_i} - \bar{Y})^2$ que de manera matricial queda como $SSE=\hat{Y}^T Y - N \bar{Y^2}$}\\
\end{enumerate}

 \item \textbf{RdeA9.1} (30 puntos) 
 Considere una ecuación para explicar los sueldos de los directores generales o CEO en términos de las ventas anuales de la empresa, el rendimiento sobre capital (\emph{roe}, en forma de porcentaje), y el rendimiento de las acciones de la empresa (\emph{ros}, en forma de porcentaje):
\[
\log(\text{salary}) \;=\; \beta_0 \;+\; \beta_1 \log(\text{sales}) \;+\; \beta_2 \,\text{roe} \;+\; \beta_3 \,\text{ros} \;+\; u.
\]


 Con los datos de \texttt{CEOSAL1}, empleando MCO se obtuvo la ecuación siguiente:
\[
\widehat{\log(\text{salary})}
= 4.32 \;+\; 0.280\,\log(\text{sales}) \;+\; 0.0174\,\text{roe} \;+\; 0.0024\,\text{ros}
\]
\[
\qquad\qquad \qquad(0.32)\qquad \qquad(0.035)\qquad\qquad(0.041)\qquad\quad (0.0054)
\]
\[
n = 209,\qquad R^2 = 0.283.
\]

\begin{enumerate}
\item (6 puntos) Establezca la hipótesis nula de que, controlando \textit{sales} y \textit{roe}, \textit{ros} no tiene efecto en el sueldo de los CEO. Establezca la alternativa de que un mejor desempeño de las acciones de la empresa incrementa el sueldo de los CEO.  Realice la prueba al nivel de significancia de 10\%.

\item (6 puntos) Establezca la hipótesis nula de que, controlando \textit{sales} y \textit{ros}, \textit{roe} no tiene efecto en el sueldo de los CEO. Establezca la alternativa de que el rendimiento sobre capital de la empresa incrementa el sueldo de los CEO.  Realice la prueba al nivel de significancia de 10\%.

\item (6 puntos) ¿Cuál es el porcentaje de aumento de \textit{salary} que se pronostica si \textit{ros} aumenta 50 puntos?  ¿Es este efecto estadísticamente significativo?

\item (6 puntos)  Realice un {test} de {significancia conjunta} de \(\{\textit{roe},\textit{ros}\}\) condicionando por \(\log(\text{sales})\), para ello considere el siguiente modelo: 


\[
\widehat{\log(\text{salary})}
= 4.82 \;+\; 0.50\,\log(\text{sales}) 
\]
\[
\qquad\qquad \qquad(0.22)\qquad \qquad(0.015)
\]
\[
n = 209,\qquad R^2 = 0.279.
\]


 ($
F_{0.90}(2,205)\approx 2.35,
\quad
F_{0.95}(2,205)\approx 3.04.
$)

\item (6 puntos)  ¿Incluiría usted \textit{roe} y \textit{ros} en el modelo final que explica las compensaciones de los CEO en términos del desempeño de la empresa? Explique.
\end{enumerate}

 \item \textbf{RdeA9.1 } (30 puntos)  Considere el siguiente modelo estimado a partir de una muestra de 120 individuos en Chile.  
El objetivo es analizar cómo la educación y el ingreso influyen en el estado de salud auto-reportado (escala 1–10):

\[
\mathrm{salud}_i = \beta_0 + \beta_1\,\mathrm{educ}_i + \beta_2\,\ln(\mathrm{ingreso}_i) + \varepsilon_i
\]

\noindent
\textbf{Descripción de variables:}
\begin{itemize}
    \item \textbf{salud}: Puntaje de salud auto-reportado (1 = muy mala, 10 = excelente).
    \item \textbf{educ}: Años de educación formal del individuo.
    \item \textbf{$\ln(\mathrm{ingreso})$}: Logaritmo natural del ingreso mensual del individuo (ingreso medido en miles de CLP).
\end{itemize}

% \noindent
% \textbf{Tabla 1: Resultados de la regresión}

\begin{table}[h!]
\centering
\begin{tabular}{lcc}
\toprule
\textbf{Variable} & \textbf{Coeficiente} & \ \\
\midrule
Intercepto         & 3.90***   \\
educación          & 0.120**   \\
$\ln(\mathrm{ingreso})$ & 0.350**  \\
\bottomrule
\end{tabular}
\vspace{0.5em}

\noindent
$R^2 = 0.36$ \\[0.3em]

\end{table}

\noindent
La matriz var–covar estimada de los coeficientes $\hat{\beta}$ es la siguiente:

\[
\widehat{\mathrm{Var}}(\hat{\beta}) =
\begin{bmatrix}
0.1600 & -0.0020 & 0.0010 \\
-0.0020 & 0.0025 & -0.0003 \\
0.0010 & -0.0003 & 0.0196
\end{bmatrix}
\]



---

\begin{flushleft}
\textbf a) (10 puntos) Interprete los signos y la significancia estadística de los coeficientes estimados para \textbf{educación} e \textbf{ingreso}.  
Explique si los resultados son consistentes con lo que se esperaría desde un punto de vista económico o social.

\vspace{0.5em}
\textbf b) (10 puntos) Testee la siguiente hipótesis: el efecto del ingreso sobre la salud es el triple del efecto de la educación sobre la salud, frente a la alternativa de que es diferente. Suponga $\alpha = 5\%$.

\vspace{0.5em}
\textbf c) (10 puntos) Pruebe conjuntamente la hipótesis de que la educación y el ingreso no influyen en el estado de salud. Suponga $\alpha = 10\%$. 

\end{flushleft}
 
    \end{enumerate}

    \section{Formulario}
\begin{itemize}

\item  
{\color{red}
$$\hat{\beta}_0 = \bar{y} - \hat{\beta}_1 \bar{x} $$
$$\hat{\beta}_1 = \frac{\sum_{i=1}^{n}(x_i - \bar{x})(y_i-\bar{y})}{\sum_{i=1}^{n}(x_i - \bar{x})^2}$$}

 {\color{red} \[ \hat{\beta} = (X^{\top}X)^{-1}X^{\top}y \]}

\item 
	{\color{red}	\[ SST = \sum_{i=1}^{n} (y_i - \bar{y})^2 \quad SSE = \sum_{i=1}^{n} (\hat{y}_i - \bar{y})^2  \quad SSR = \sum_{i=1}^{n} \hat{u}_i^2\]}
		

    \item {\color{red} \[ R^2  = 1 - \frac{SSR}{SST} \quad \quad \bar{R}^2 =  1 - \frac{(n-1)SSR}{(n-k-1)SST} \]}

 \item {\color{red} \[ t_{\hat{\beta_j}}=\frac{\hat{\beta}_j - \beta_j}{\text{s.e.}(\hat{\beta}_j)} \sim t_{n-k-1} \]}

 \item {\color{red}\[
    F = \frac{SSR_{r} - SSR_{ur}}{SSR_{ur}} \cdot \frac{n - k - 1}{j} \sim F(j, n - k - 1)
    \]}

\item {\color{red}%
\[ F_{0.90}(2,117) \approx 2.30 \]%
}

    \begin{figure}[H]
    \centering
    \includegraphics[width=1.0\linewidth]{certamenes/norm_bilat_ren.png}
\end{figure}
\end{itemize}   

    \end{document}